% Recovered verbatim from chapters_31_35_with_part/chapters_31_35_with_part.tex
% Canonical chapter 32
\chapter{Chromatic Distributions and Structure Formation}

\section{Chromatic Seeds for Density Perturbations}

Small spatial variations in $\chi_a(x)$ produce fluctuations in the effective energy density:
\[
  \delta \rho_\chi \approx G^{ab}\dot{\chi}_{a,0}\,\delta\dot{\chi}_b + \partial_a U(\bm{\chi}_0)\,\delta\chi_a,
\]
which act as seeds for density perturbations.
These seeds are processed by the cosmological expansion and gravitational dynamics, leading to the growth of structure.

\section{Power Spectra and Deviations from \texorpdfstring{$\Lambda$}{Lambda}CDM}

The primordial power spectrum $P(k)$ can be affected by:
\begin{itemize}
  \item the initial chromatic configuration $\bm{\chi}_0$,
  \item the shape of the chromatic potential $U(\bm{\chi})$,
  \item couplings to the phase field and ER network.
\end{itemize}
This may lead to small deviations from the nearly scale-invariant spectra of standard $\Lambda$CDM, potentially visible as features or running in $P(k)$.

\section{Non-Gaussianity and ER-Induced Correlations}

Nonlinearities in chromatic dynamics and ER connectivity can generate non-Gaussian correlations in the matter distribution.
These appear as higher-order moments of the density field and may be probed through large-scale structure surveys and CMB observations.

